\documentclass[10.5pt]{article}
\usepackage[letterpaper, margin=0.6in, top=0.5in, bottom=0.5in]{geometry}
\usepackage{amsmath,amssymb,amsthm}
\usepackage{fancyhdr}
\usepackage{xcolor}
\usepackage{tikz}
\usepackage{enumitem}
\usepackage{tcolorbox}
\usepackage{hyperref}

% Define colors matching instructor board annotations
\definecolor{chalkblue}{RGB}{0, 105, 185}
\definecolor{chalkred}{RGB}{205, 35, 35}
\definecolor{chalkpink}{RGB}{230, 90, 120}
\definecolor{chalkcyan}{RGB}{15, 160, 190}
\definecolor{chalkgreen}{RGB}{20, 130, 50}
\definecolor{boxbg}{RGB}{252, 252, 254}
\definecolor{boxframe}{RGB}{150, 160, 180}

\pagestyle{fancy}
\fancyhf{}
\renewcommand{\headrulewidth}{0pt}
\renewcommand{\footrulewidth}{0pt}

\setlength{\parindent}{0pt}
\setlength{\parskip}{2.5pt}

\begin{document}

% ==============================================================================
% PAGE 1
% ==============================================================================
\noindent \textbf{\large Math 105L} \hfill \textbf{\large Lecture 4-2}\\[1pt]
{\textbf{\Large Rates of Change: From Average Velocity to Instantaneous Velocity}}\\[-5pt]
\rule{\textwidth}{1.2pt}

\vspace{0.03in}

\textbf{\large 1. The Geometry of Average Velocity}

The formula for average velocity over a time interval $[a, b]$ has the exact mathematical form of a line slope:
\[
  \text{Average Velocity on } [a, b] = \frac{\Delta s}{\Delta t} = \frac{s(b) - s(a)}{b - a}
\]
Setting $b = a + h$, the interval is $[a, a + h]$ and the average velocity becomes the \textbf{difference quotient}:
\[
  \mathbf{\text{Avg Vel on } [a, a+h] = \frac{s(a+h) - s(a)}{h}}
\]
\textbf{Geometric Interpretation:} The average velocity equals the slope of the \textbf{secant line} connecting $(a, s(a))$ and $(a+h, s(a+h))$.

\vspace{0.03in}

\textbf{\large 2. Classroom Case Study: Projectile Motion $s(t) = -4t(t-1)$}

Consider an object moving along a line with position $s(t) = -4t(t-1) = -4t^2 + 4t$ (in meters) from $t = 0$ to $t = 1$ seconds:

\begin{itemize}[leftmargin=1.2em, topsep=1pt, itemsep=3pt]
  \item \textbf{Interval $[0, 1]$:}
  \[
    \text{Avg Vel} = \frac{s(1) - s(0)}{1 - 0} = \frac{0 - 0}{1} = \mathbf{0 \text{ m/s}} \quad \text{(secant line is horizontal)}
  \]
  \item \textbf{Interval $[0, 0.5]$:}
  \[
    \text{Avg Vel} = \frac{s(0.5) - s(0)}{0.5 - 0} = \frac{1 - 0}{0.5} = \mathbf{2 \text{ m/s}}
  \]
  \item \textbf{Interval $[0, 0.1]$:}
  \[
    s(0.1) = -4(0.1)(-0.9) = 0.36 \implies \text{Avg Vel} = \frac{0.36 - 0}{0.1} = \mathbf{3.6 \text{ m/s}}
  \]
  \item \textbf{Interval $[0, 0.01]$:}
  \[
    s(0.01) = -4(0.01)(-0.99) = 0.0396 \implies \text{Avg Vel} = \frac{0.0396 - 0}{0.01} = \mathbf{3.96 \text{ m/s}}
  \]
\end{itemize}

\vspace{0.03in}

\begin{tcolorbox}[colback=boxbg, colframe=boxframe, boxrule=0.8pt, arc=3pt, left=8pt, right=8pt, top=4pt, bottom=4pt]
\textbf{\large The Instantaneous Velocity at $t = 0$}
As the time interval $h$ shrinks toward zero, the secant lines pivot toward the \textbf{tangent line} at $(0,0)$:
\[
  v(0) = \lim_{h\to 0} \frac{s(0+h) - s(0)}{h} = \lim_{h\to 0}\frac{-4h(h-1)}{h} = \lim_{h\to 0}\big[-4(h-1)\big] = \mathbf{4 \text{ m/s}}
\]
\end{tcolorbox}

\vspace{0.04in}

\textbf{\large 3. Formal Summary}

\begin{tcolorbox}[colback=chalkblue!6, colframe=chalkblue, boxrule=1pt, arc=3pt, left=8pt, right=8pt, top=4pt, bottom=4pt]
\begin{itemize}[leftmargin=1.2em, topsep=1pt, itemsep=3pt]
  \item \textbf{Instantaneous Velocity Definition:}
  \[
    \mathbf{v(a) = \lim_{h\to 0} \frac{s(a+h) - s(a)}{h}}
  \]
  \item \textbf{Graphical Signature:} On the graph of position $s(t)$, the instantaneous velocity at $t = a$ is given by the \textbf{slope of the tangent line} to the curve at the point $(a, s(a))$.
\end{itemize}
\end{tcolorbox}

\end{document}

